\daytitle{23}{电磁感应}

\begin{learningbox}{本节导读}
前两节研究“电流产生磁场”和“磁场对电流、运动电荷的作用”。本节反过来看：穿过回路的磁通量发生变化时，会产生怎样的电磁效应？这就是电磁感应。接下来学习磁通量、感应电动势和感应电流的产生条件、楞次定律、法拉第电磁感应定律，以及导体棒切割磁感线和线圈进出磁场的典型模型。
\end{learningbox}

\textbf{本节读题路线：} 先问“磁通量是否变化”，决定有没有感应电动势；再用楞次定律问“感应电流要产生怎样的磁场”，判断方向；最后用法拉第定律问“变化有多快”，计算大小。这样处理，电磁感应题就不再是零散方向规则的记忆，而是一条因果链。

\textbf{符号提醒：} 本节统一用 $\mathcal E$ 表示\textbf{感应电动势}，用 $E$ 表示电场强度；用 $Q$ 表示某过程中通过导体横截面的电荷量。

\section*{一、学习目标}
\begin{enumerate}
  \item 能由 $\Phi=\vec B\cdot\vec S=BS\cos\theta$ 判断磁通量的大小和正负。
  \item 能区分产生感应电动势与形成感应电流的条件。
  \item 会用楞次定律判断感应电流方向。
  \item 能区分平均感应电动势 $\overline{\mathcal E}$ 与瞬时感应电动势 $\mathcal E$。
  \item 能说明电阻不变时净通过电荷量与磁通量变化的关系及其适用条件。
  \item 会分析导体棒切割磁感线、线圈进出磁场和发电机的基本模型。
\end{enumerate}
\section*{二、磁通量变化怎样产生电流}

初中知道发电机能把机械能转化为电能，通常先从现象入手：线圈在磁场中转动能发电。高中要回答三个更精确的问题：
\begin{itemize}
  \item 什么时候会产生感应电流？
  \item 感应电流方向怎么判断？
  \item 感应电动势大小由什么决定？
\end{itemize}

第一个问题要分成两层：磁通量变化会产生感应电动势，导电回路闭合后才形成感应电流。后两个问题分别由楞次定律和法拉第电磁感应定律回答。

\textbf{科学发展视角：} 电磁感应连接了发电、输电和用电。发电机、无线充电、电磁炉、金属探测和电动汽车的再生制动，都建立在“磁通量变化产生感应电动势”的思想上。学到这里，电和磁不再是两条线，而开始合成一个完整的电磁学框架。

\section*{三、法拉第定律与楞次定律}

\begin{conceptbox}{磁通量}

磁通量描述磁场穿过某个有向面积的程度。先为回路所围面积选定一个正法线方向，并定义面积矢量 $\vec S$：方向沿正法线，大小等于面积 $S$。若磁场为匀强磁场，则
\[
\boxed{\Phi=\vec B\cdot\vec S=BS\cos\theta}
\]
其中 $\theta$ 是 $\vec B$ 与正法线（即 $\vec S$）的夹角。正法线一经选定，磁通量就有正负；若只问穿过面积的磁通量大小，再取 $|\Phi|$。

\begin{itemize}
  \item 单位：韦伯（$\mathrm{Wb}$）。
  \item 若磁场沿正法线垂直穿过平面，$\theta=0^\circ$，$\Phi=BS$。
  \item 若磁场沿正法线反方向垂直穿过平面，$\theta=180^\circ$，$\Phi=-BS$。
  \item 若磁场平行于平面，$\theta=90^\circ$，$\Phi=0$。
\end{itemize}

\begin{center}
% TikZ 图：磁通量角度
% Day 23：电磁感应
\begin{tikzpicture}[scale=0.9,>=Latex]
  \foreach \x in {-2.2,-1.5,-0.8,-0.1,0.6,1.3,2.0} {
    \draw[->,blue,thick] (\x,-1.45) -- (\x,1.45);
  }
  \node[blue] at (2.55,1.35) {$B$};
  \draw[thick,fill=gray!12] (-1.2,-0.55) -- (1.2,-0.25) -- (1.55,0.75) -- (-0.85,0.45) -- cycle;
  \draw[->,red,thick] (0.15,0.1) -- (0.75,1.15);
  \node[red,fill=white,inner sep=1pt] at (-0.08,1.28) {法线};
  \draw[->,orange,thick] (0.15,0.1) -- (0.15,1.2);
  \draw[orange] (0.15,0.65) arc (90:60:0.55);
  \node[orange] at (0.48,0.78) {$\theta$};
  \node at (0,-1.9) {$\Phi=BS\cos\theta$};
\end{tikzpicture}


\figcaption{选定正法线后，磁通量由磁场强弱、面积大小和面积朝向共同决定}\label{fig:magnetic-flux-angle}
\end{center}

\end{conceptbox}

\begin{conceptbox}{电磁感应现象}

当穿过回路所围面积的磁通量发生变化时，回路中会产生感应电动势；若导电回路闭合，才会形成感应电流。计算普通电阻回路中的电流时，还要知道回路的总电阻。

\[
\boxed{\text{磁通量变化}\Rightarrow\text{感应电动势}\quad;\quad
  \text{导电回路闭合}\Rightarrow\text{形成感应电流}}
\]

磁通量可以通过三种方式改变：
\begin{itemize}
  \item 改变磁场强弱 $B$；
  \item 改变回路面积 $S$；
  \item 改变回路与磁场的夹角 $\theta$。
\end{itemize}

\begin{center}
% TikZ 图：改变磁通量的三种方式
% Day 23：电磁感应
\begin{tikzpicture}[scale=0.82,>=Latex]
  \draw[thick] (-3.8,-0.75) rectangle (-2.4,0.75);
  \foreach \x in {-3.5,-3.1,-2.7} {
    \draw[->,blue] (\x,-0.55) -- (\x,0.55);
  }
  \node at (-3.1,-1.15) {改变 $B$};
  \draw[thick] (-0.7,-0.75) rectangle (0.7,0.75);
  \draw[dashed,thick] (-1.0,-1.0) rectangle (1.0,1.0);
  \foreach \x in {-0.5,0,0.5} {
    \draw[->,blue] (\x,-1.05) -- (\x,1.05);
  }
  \node at (0,-1.35) {改变 $S$};
  \draw[thick,fill=gray!12] (2.35,-0.55) -- (3.65,-0.25) -- (3.9,0.6) -- (2.6,0.3) -- cycle;
  \draw[thick,dashed] (2.35,-0.75) rectangle (3.85,0.75);
  \foreach \x in {2.55,3.05,3.55} {
    \draw[->,blue] (\x,-0.95) -- (\x,0.95);
  }
  \node at (3.1,-1.35) {改变角度};
\end{tikzpicture}


\figcaption{改变 $B$、$S$ 或角度，都可能引起磁通量变化}\label{fig:induction-three-ways}
\end{center}

\textbf{边界提醒：} 如果回路不闭合，仍可能在断口两端出现感应电动势，但不能形成沿整个回路持续流动的感应电流。

\end{conceptbox}

\begin{conceptbox}{楞次定律}

闭合回路中，感应电流产生的磁场总要阻碍引起感应电流的磁通量变化。

\[
\boxed{\text{感应电流的磁场，总要阻碍引起它的磁通量变化}}
\]

\begin{center}
% TikZ 图：楞次定律
% Day 23：电磁感应
\begin{tikzpicture}[scale=0.9,>=Latex]
  \draw[thick,fill=gray!12] (-3.0,-0.45) rectangle (-1.4,0.45);
  \node at (-2.6,0) {$N$};
  \node at (-1.8,0) {$S$};
  \draw[->,red,thick] (-1.25,0) -- (-0.55,0) node[midway,above] {靠近};
  \draw[thick] (0.15,0) ellipse (0.35 and 1.05);
  \draw[thick] (0.45,0) ellipse (0.35 and 1.05);
  \draw[thick] (0.75,0) ellipse (0.35 and 1.05);
  \draw[->,blue,thick] (-0.55,0.72) -- (-1.25,0.72) node[left] {阻碍};
  \node at (0.45,-1.45) {线圈近端表现为 $S$ 极，阻碍磁铁靠近};
\end{tikzpicture}


\figcaption{磁铁靠近线圈时，线圈产生的磁场阻碍磁通量增加}\label{fig:lenz-law-magnet}
\end{center}

\textbf{四步判断法：}
\begin{enumerate}
  \item 判断原磁场方向；
  \item 判断穿过回路的磁通量是增加还是减少；
  \item 用“阻碍变化”确定感应磁场方向；
  \item 用右手螺旋定则确定感应电流方向。
\end{enumerate}

\textbf{方向判断清单：} 在原磁场方向不变、只比较磁通量大小增减的常见模型中，
\[
\begin{array}{c|c}
\text{原磁通量变化} & \text{感应磁场方向} \\
\hline
\text{原磁通量增加} & \text{与原磁场方向相反，阻碍增加} \\
\text{原磁通量减少} & \text{与原磁场方向相同，阻碍减少}
\end{array}
\]

\end{conceptbox}

\begin{conceptbox}{法拉第电磁感应定律}

设线圈有 $N$ 匝，$\Phi$ 表示每一匝的磁通量。在一段时间 $\Delta t$ 内，平均感应电动势为
\[
\boxed{\overline{\mathcal E}=-N\frac{\Delta\Phi}{\Delta t}},
\qquad
\boxed{\left|\overline{\mathcal E}\right|=N\left|\frac{\Delta\Phi}{\Delta t}\right|}.
\]
负号表示感应效应遵循楞次定律。若磁通量随时间连续变化，某一时刻的瞬时感应电动势为
\[
\boxed{\mathcal E=-N\frac{\mathrm d\Phi}{\mathrm dt}}.
\]

\begin{itemize}
  \item $N$ 是线圈匝数。
  \item $\Delta\Phi$ 是每一匝线圈的磁通量变化量。
  \item 平均变化率决定平均感应电动势，瞬时变化率决定瞬时感应电动势。
\end{itemize}

\begin{center}
\begin{tikzpicture}[scale=0.9,>=Latex,line cap=round,line join=round]
  \path[use as bounding box] (-0.55,-0.75) rectangle (6.45,3.35);
  \draw[->,thick] (0,0) -- (6.0,0) node[right] {$t$};
  \draw[->,thick] (0,0) -- (0,3.0) node[above] {$\Phi$};

  \draw[blue,very thick]
    (0.45,0.38) .. controls (1.35,0.55) and (2.20,1.18) ..
    (3.05,1.72) .. controls (4.15,2.42) and (5.05,2.58) .. (5.55,2.68);

  \coordinate (Pone) at (1.15,0.67);
  \coordinate (Ptwo) at (4.75,2.51);
  \coordinate (P) at (3.05,1.72);
  \fill (Pone) circle (1.7pt) node[below right,font=\small] {$P_1$};
  \fill (Ptwo) circle (1.7pt) node[above left,font=\small] {$P_2$};
  \fill[orange!85!black] (P) circle (1.9pt) node[below right,font=\small] {$P$};

  \draw[dashed,gray] (Pone) -- (1.15,0);
  \draw[dashed,gray] (Ptwo) -- (4.75,0);
  \draw[dashed,red,thick] (Pone) -- (Ptwo);
  \node[red,font=\small,fill=white,inner sep=1pt] at (2.55,1.02)
    {割线斜率 $\dfrac{\Delta\Phi}{\Delta t}$};

  \draw[orange!85!black,thick] (1.95,1.00) -- (4.12,2.36);
  \node[orange!85!black,font=\small,fill=white,inner sep=1pt] at (4.42,1.94)
    {切线斜率 $\dfrac{\mathrm d\Phi}{\mathrm dt}$};

  \node[font=\small] at (3.0,-0.52)
    {$\left|\overline{\mathcal E}\right|=N\left|\dfrac{\Delta\Phi}{\Delta t}\right|,
      \quad |\mathcal E|=N\left|\dfrac{\mathrm d\Phi}{\mathrm dt}\right|$};
\end{tikzpicture}


\figcaption{$\Phi$-$t$ 图像的割线斜率对应平均感应电动势，切线斜率对应瞬时感应电动势}\label{fig:faraday-flux-time}
\end{center}

\textbf{认知提醒：} 图像斜率为负不等于“感应电动势一定为负”。由 $\mathcal E=-N\,\mathrm d\Phi/\mathrm dt$ 可知，感应电动势的正负还取决于预先选定的回路正方向；只比较大小时取变化率的绝对值。

\end{conceptbox}

\begin{samepage}
\section*{四、感应电动势的关系式}

\textbf{列式前检查：}
\nopagebreak[4]
\begin{enumerate}
  \item 先选定回路、正法线和与之匹配的回路正方向，再判断磁通量 $\Phi$ 是否变化。
  \item 用楞次定律时，先说“原磁通怎样变”，再说“感应磁场怎样阻碍这种变化”。
  \item 用 $\left|\overline{\mathcal E}\right|=N|\Delta\Phi/\Delta t|$ 前先确认求平均电动势；若要求某一时刻大小，要看瞬时变化率或导体棒切割模型。
  \item 导体棒切割磁感线时，要判断速度方向、磁场方向和有效长度。
  \item 涉及能量时，外力做功通常转化为电能和焦耳热，不能只看电流方向。
\end{enumerate}
\end{samepage}

\begin{center}
\begin{tabular}{lll}
\toprule
公式 & 物理意义 & 适用条件 \\
\midrule
$\Phi=\vec B\cdot\vec S$ & 有向磁通量 & 先选正法线；匀强磁场、平面面积 \\
$\left|\overline{\mathcal E}\right|=N\left|\dfrac{\Delta\Phi}{\Delta t}\right|$ & 平均感应电动势 & 一段时间内的平均值 \\
$\mathcal E=-N\dfrac{\mathrm d\Phi}{\mathrm dt}$ & 瞬时感应电动势 & 磁通量连续变化 \\
$|Q_{\text{净}}|=\dfrac{N|\Delta\Phi|}{R}$ & 净通过电荷量 & $R$ 不变，忽略自感 \\
$\mathcal E=BLv$ & 导体棒切割磁感线 & $B$、$L$、$v$ 两两垂直 \\
$I=\dfrac{\mathcal E}{R}$ & 感应电流 & 闭合电阻回路总电阻为 $R$ \\
$F_A=BIL$ & 安培力（阻碍棒运动） & 导体棒电流与匀强磁场垂直 \\
\bottomrule
\end{tabular}
\end{center}

\begin{learningbox}{电磁感应题的四问分析清单}
电磁感应题不要从“电流方向”开始猜。先按四个问题排顺序：
\begin{enumerate}
  \item \textbf{有没有：} 选定回路，判断穿过它的磁通量 $\Phi$ 是否变化；不变则无感应电动势。
  \item \textbf{方向：} 若回路闭合，用楞次定律先定感应磁场方向，再用右手螺旋定则定感应电流方向。
  \item \textbf{大小：} 平均值用 $\left|\overline{\mathcal E}\right|=N\left|\dfrac{\Delta\Phi}{\Delta t}\right|$；瞬时值看 $\Phi$-$t$ 图像切线斜率；导体棒垂直切割可用 $\mathcal E=BLv$。
  \item \textbf{能量：} 若有外力拉动、线圈进出磁场或导体棒运动，要追问机械能、电能和焦耳热之间如何转化。
\end{enumerate}
方向判断和大小计算要分开：楞次定律负责判断方向，法拉第定律负责计算大小。
\end{learningbox}

\section*{五、线圈与导体棒模型}

\subsection*{模型 1：导体棒切割磁感线}

先看导体棒与两条静止导轨构成的闭合回路。若匀强磁场 $B$ 垂直导轨平面，导体棒的有效长度 $L$ 与速度 $v$ 垂直，则棒运动时回路面积发生变化，感应电动势为
\[
\boxed{\mathcal E=BLv}.
\]

这个公式可以由法拉第电磁感应定律推出。导体棒在时间 $\Delta t$ 内前进 $\Delta x$，回路面积增加
\[
\Delta S=L\Delta x.
\]
匀强磁场垂直穿过回路，磁通量变化
\[
\Delta\Phi=B\Delta S=BL\Delta x,
\]
所以
\[
\mathcal E=\left|\frac{\Delta\Phi}{\Delta t}\right|=BL\frac{\Delta x}{\Delta t}=BLv.
\]
因此，$\mathcal E=BLv$ 不是孤立的新公式，而是法拉第电磁感应定律在这一闭合导轨模型中的具体结果。对于单独运动的导体棒，也可以用自由电荷所受洛伦兹力解释棒两端的电势差。

\textbf{适用提醒：} 若导体棒不是完全垂直切割，或只有一部分长度处在磁场中，要先取有效长度、有效速度分量，或直接回到磁通量变化率，不能直接套用整根棒的 $BLv$。

\begin{center}
% TikZ 图：导体棒切割磁感线
% Day 23：电磁感应
\begin{tikzpicture}[scale=0.9,>=Latex]
  \draw[thick] (-2,1) -- (2.4,1);
  \draw[thick] (-2,-1) -- (2.4,-1);
  \draw[thick] (-2,1) -- (-2,-1);
  \draw[very thick,red] (0.1,-1) -- (0.1,1) node[midway,right] {$L$};
  \draw[->,red,very thick] (0.1,0) -- (1.25,0) node[right] {$v$};
  \foreach \x in {-1.5,-0.8,-0.1,0.6,1.3,2.0} {
    \foreach \y in {-0.55,0.25} {
      \draw[blue,thick] (\x-0.07,\y-0.07) -- (\x+0.07,\y+0.07);
      \draw[blue,thick] (\x-0.07,\y+0.07) -- (\x+0.07,\y-0.07);
    }
  }
  \node[blue] at (2.85,0.85) {$B$ 向里};
  \draw[->,orange,thick] (-0.6,1.2) -- (-1.55,1.2) node[midway,above] {$I$};
  \draw[->,orange,thick] (0.32,-0.65) -- (0.32,0.65);
  \node at (0,-1.55) {$E=BLv,\quad I=E/R$};
\end{tikzpicture}


\figcaption{导体棒切割磁感线：面积增大，磁通量增大，产生感应电流}\label{fig:motional-emf-rail}
\end{center}

若回路总电阻为 $R$，则
\[
I=\frac{BLv}{R}
\]
闭合回路形成感应电流后，导体棒还会受到安培力。按楞次定律判断，它的效果是阻碍磁通量的变化；在这里的单棒切割模型中，通常表现为阻碍导体棒相对磁场的运动。

\subsection*{模型 2：线圈进入或离开磁场}

设有界匀强磁场区域的宽度大于线圈沿运动方向的宽度。线圈进入磁场时，处在磁场中的有效面积增大；完全进入后，有效面积不变；离开磁场时，有效面积减小。

\begin{center}
% TikZ 图：线圈进入磁场
% Day 23：电磁感应
\begin{tikzpicture}[scale=0.88,>=Latex]
  \draw[thick,fill=blue!4] (-0.2,-1.15) rectangle (3.5,1.15);
  \node[blue] at (2.8,0.92) {匀强磁场};
  \foreach \x in {0.3,1.0,1.7,2.4,3.1} {
    \foreach \y in {-0.75,0,0.75} {
      \draw[blue,thick] (\x-0.055,\y-0.055) -- (\x+0.055,\y+0.055);
      \draw[blue,thick] (\x-0.055,\y+0.055) -- (\x+0.055,\y-0.055);
    }
  }
  \draw[thick,red] (-1.45,-0.55) rectangle (-0.45,0.55);
  \draw[->,red,thick] (-1.0,0) -- (-0.25,0) node[above] {$v$};
  \node[align=center] at (-1.0,-1.75) {进入\\$\Phi$ 增大};
  \draw[thick,orange] (1.0,-0.55) rectangle (2.0,0.55);
  \draw[->,orange,thick] (1.35,0) -- (2.1,0);
  \node[align=center] at (1.5,-1.75) {内部\\$\Phi$ 不变};
  \draw[thick,purple] (3.0,-0.55) rectangle (4.0,0.55);
  \draw[->,purple,thick] (3.35,0) -- (4.1,0);
  \node[align=center] at (3.5,-1.75) {离开\\$\Phi$ 减小};
\end{tikzpicture}


\figcaption{线圈进出磁场：只有磁通量变化阶段才有感应电流}\label{fig:coil-enter-field}
\end{center}

如果矩形线圈匀速穿过边界清晰的匀强磁场，进入和离开阶段的有效面积随时间线性变化，因此 $\Phi$-$t$ 图像是“上升—水平—下降”的分段图像。感应电动势看斜率：若进、出磁场时的磁场强度、速度和有效边长分别相同，则两阶段的感应电动势大小相同、方向相反，与回路电阻无关；若还要比较感应电流大小，才需要比较回路电阻。完全处在匀强磁场中且面积、取向不变时，感应电动势为零。

\begin{center}
% TikZ 图：线圈进出磁场的磁通量与感应电动势图像
% Day 23：电磁感应
\begin{tikzpicture}[scale=0.9,>=Latex]
  \path[use as bounding box] (-0.35,-2.45) rectangle (6.45,2.35);

  % Magnetic flux graph.
  \draw[->] (0,0) -- (6.05,0) node[right] {$t$};
  \draw[->] (0,0) -- (0,2.05) node[above] {$\Phi$};
  \draw[blue,very thick] (0.45,0) -- (2.05,1.45) -- (3.85,1.45) -- (5.45,0);
  \draw[dashed] (2.05,1.45) -- (2.05,-1.55);
  \draw[dashed] (3.85,1.45) -- (3.85,-1.55);
  \node[font=\small] at (1.25,1.85) {进入};
  \node[font=\small] at (2.95,1.85) {内部};
  \node[font=\small] at (4.65,1.85) {离开};

  % Induced emf/current sign graph.
  \draw[->] (0,-1.65) -- (6.05,-1.65) node[right] {$t$};
  \draw[->] (0,-2.25) -- (0,-0.85) node[above] {$\mathcal{E}$};
  \draw[red,very thick] (0.45,-2.25) -- (2.05,-2.25);
  \draw[red,very thick] (2.05,-1.65) -- (3.85,-1.65);
  \draw[red,very thick] (3.85,-1.00) -- (5.45,-1.00);
  \draw[red,dashed] (2.05,-2.25) -- (2.05,-1.65);
  \draw[red,dashed] (3.85,-1.65) -- (3.85,-1.00);
  \node[red,font=\small] at (1.10,-2.42) {$-\mathcal E_0$};
  \node[red,font=\small] at (2.95,-1.38) {$0$};
  \node[red,font=\small] at (4.70,-0.58) {$+\mathcal E_0$};

\end{tikzpicture}


\figcaption{线圈进出磁场的分段图像：看斜率判断感应电动势}\label{fig:coil-flux-current-graph}
\end{center}

\subsection*{模型 3：感应电荷量模型}

有些题不问瞬时电流或最大电流，而问“整个过程中通过导体横截面的电荷量”。先选定回路正方向，并假定总电阻 $R$ 保持不变、线圈自感可以忽略。由
\[
\mathcal E=-N\frac{\mathrm d\Phi}{\mathrm dt},
\qquad I=\frac{\mathcal E}{R},
\]
可得带符号的净通过电荷量
\[
\boxed{Q_{\text{净}}=\int_{t_1}^{t_2}I\,\mathrm dt
=\frac{N(\Phi_1-\Phi_2)}{R}}.
\]
因此
\[
\boxed{|Q_{\text{净}}|=\frac{N|\Delta\Phi|}{R}}.
\]
若磁通量单调变化、感应电流方向始终不变，$|Q_{\text{净}}|$ 就等于整个过程中累计通过横截面的电荷量。若电流中途反向，应分段计算并区分“各段电荷量绝对值之和”与“带符号的净电荷量”。

\begin{center}
% TikZ 图：感应电荷量与 I-t 图像面积
% Day 23：电磁感应
\begin{tikzpicture}[scale=0.92,>=Latex]
  \tikzset{
    axis/.style={->,thick},
    curve/.style={very thick,blue},
    shadearea/.style={fill=blue!18,draw=none},
    note/.style={font=\small,align=center}
  }

  \begin{scope}[xshift=-3.1cm]
    \draw[axis] (0,0) -- (2.8,0) node[right] {$t$};
    \draw[axis] (0,0) -- (0,2.1) node[above] {$I$};
    \path[shadearea] (0.35,0) -- (0.35,1.7)
      .. controls (0.65,1.95) and (1.05,1.95) .. (1.35,1.7)
      -- (1.35,0) -- cycle;
    \draw[curve] (0.35,0) -- (0.35,1.7)
      .. controls (0.65,1.95) and (1.05,1.95) .. (1.35,1.7)
      -- (1.35,0);
    \node[note] at (1.35,-0.48) {变化快\\电流大、时间短};
    \node[blue,font=\small] at (1.85,1.25) {$q=$面积};
  \end{scope}

  \begin{scope}[xshift=3.1cm]
    \draw[axis] (0,0) -- (2.8,0) node[right] {$t$};
    \draw[axis] (0,0) -- (0,2.1) node[above] {$I$};
    \path[shadearea] (0.35,0) -- (0.35,0.9)
      .. controls (0.95,1.10) and (1.95,1.10) .. (2.55,0.9)
      -- (2.55,0) -- cycle;
    \draw[curve] (0.35,0) -- (0.35,0.9)
      .. controls (0.95,1.10) and (1.95,1.10) .. (2.55,0.9)
      -- (2.55,0);
    \node[note] at (1.45,-0.48) {变化慢\\电流小、时间长};
    \node[blue,font=\small] at (1.85,0.63) {$q=$面积};
  \end{scope}

  \node[note] at (0,2.55) {若 $\Delta\Phi$ 相同且 $R$ 不变，两图面积相同};
  \node[note] at (0,-1.18) {$q=\bar I\Delta t=\dfrac{n|\Delta\Phi|}{R}$};
\end{tikzpicture}


\figcaption{磁通量单调变化时，$I$-$t$ 图线与时间轴围成的面积表示通过横截面的电荷量}\label{fig:induction-charge-area}
\end{center}

\textbf{认知提醒：} 对于磁通量单调变化的两个过程，变化快时通常电流大、时间短，变化慢时通常电流小、时间长；若 $N$、$R$ 和 $|\Delta\Phi|$ 分别相同，则累计通过的电荷量相同。这个结论不能用来判断瞬时电流大小，也不适用于电阻明显变化或电流多次反向而未分段处理的过程。

\subsection*{模型 4：发电机模型}

矩形线圈绕位于线圈平面内的固定轴，在匀强磁场中以角速度 $\omega$ 匀速转动。若从线圈面积法线与磁场方向重合的位置开始计时，则
\[
\Phi=BS\cos\omega t.
\]
对 $N$ 匝线圈，瞬时感应电动势为
\[
\mathcal E=-N\frac{\mathrm d\Phi}{\mathrm dt}
=NBS\omega\sin\omega t
=\mathcal E_m\sin\omega t,
\qquad
\mathcal E_m=NBS\omega.
\]
可见，机械转动使面积法线与磁场的夹角周期性变化，从而产生交变电动势；其周期由转动角速度决定，最大值由 $N$、$B$、$S$ 和 $\omega$ 共同决定。

\begin{center}
\begin{tikzpicture}[scale=0.92,>=Latex,line cap=round,line join=round]
  \path[use as bounding box] (-3.45,-2.25) rectangle (3.45,2.35);

  \draw[thick,fill=gray!12] (-3.15,-1.15) rectangle (-2.45,1.15);
  \node at (-2.80,0) {$N$};
  \draw[thick,fill=gray!12] (2.45,-1.15) rectangle (3.15,1.15);
  \node at (2.80,0) {$S$};
  \foreach \y in {-0.78,-0.26,0.26,0.78} {
    \draw[->,blue,thick] (-2.35,\y) -- (2.35,\y);
  }
  \node[blue,fill=white,inner sep=1pt] at (-1.82,1.02) {$\vec B$};

  % A wire loop seen obliquely. Its vertical axis lies in the loop plane.
  \coordinate (T) at (-0.12,1.25);
  \coordinate (B) at (-0.12,-1.25);
  \coordinate (LT) at (-1.05,0.82);
  \coordinate (LB) at (-1.05,-0.82);
  \coordinate (RT) at (0.82,0.66);
  \coordinate (RB) at (0.82,-0.66);
  \draw[very thick] (LT) -- (RT) -- (RB) -- (LB) -- cycle;
  \draw[dashed,gray,thick] (T) -- (B);
  \fill[gray] (T) circle (1.7pt);
  \fill[gray] (B) circle (1.7pt);
  \node[gray!70!black,font=\small,anchor=west] at (0.02,-1.35) {转轴};

  \draw[->,orange!85!black,very thick] (-0.10,0) -- (1.15,0.72)
    node[above right,font=\small] {$\vec S$};
  \draw[orange!85!black,thick] (0.46,0) arc (0:30:0.56);
  \node[orange!85!black,font=\small] at (0.70,0.19) {$\theta$};

  \draw[->,red,thick] (0.46,1.43) arc (15:330:0.58 and 0.24);
  \node[red,font=\small] at (0.48,1.86) {$\omega$};

  \node[font=\small] at (0,-1.93)
    {$\Phi=BS\cos\theta,\qquad \theta=\omega t$};
\end{tikzpicture}


\figcaption{发电机基本模型：线圈绕平面内的轴转动，面积法线与磁场的夹角周期性变化}\label{fig:generator-model}
\end{center}

\section*{六、例题精讲}

\begin{examplebox}{磁通量的计算}
\textbf{题意分析：} 面积 $S=0.20\ \mathrm{m^2}$ 的线圈放在 $B=0.50\ \mathrm{T}$ 的匀强磁场中。若磁场方向与线圈平面法线夹角为 $60^\circ$，求磁通量。

\textbf{物理模型：} 磁通量 $\Phi=BS\cos\theta$。

\textbf{解题过程：}
\[
\Phi=BS\cos60^\circ
=0.50\times0.20\times\frac12
=0.050\ \mathrm{Wb}
\]

\textbf{方法总结：} 磁通量题先画平面法线；$\Phi=BS\cos\theta$ 中 $\theta$ 是磁场与法线夹角，不是与平面夹角。
\end{examplebox}

\begin{examplebox}{用楞次定律判断感应电流方向}
\textbf{题意分析：} 一个闭合圆形线圈固定在纸面内，垂直纸面向里的匀强磁场正在增强。判断线圈中感应磁场和感应电流的方向。

\textbf{物理模型：} 原磁场方向不变，向里的磁通量正在增加。依次判断“原磁通量怎样变—感应磁场怎样阻碍—感应电流怎样绕”。

\begin{center}
% TikZ 图：楞次定律方向判断例题
% Day 23：电磁感应
\begin{tikzpicture}[scale=0.92,>=Latex,line cap=round,line join=round]

  \foreach \x in {-2.35,-1.55,-0.75,0.75,1.55,2.35} {
    \foreach \y in {-0.95,0,0.95} {
      \draw[blue,thick] (\x-0.07,\y-0.07) -- (\x+0.07,\y+0.07);
      \draw[blue,thick] (\x-0.07,\y+0.07) -- (\x+0.07,\y-0.07);
    }
  }
  \node[blue,fill=white,inner sep=1pt] at (-2.25,1.55)
    {原磁场向里且增强};

  \draw[very thick] (0,0) circle (1.35);
  \draw[->,red,very thick] (12:1.35) arc (12:305:1.35);
  \node[red,fill=white,inner sep=1pt] at (1.85,-1.30)
    {$I_{\text{感}}$ 逆时针};

  \fill[orange!85!black] (0,0) circle (2.4pt);
  \draw[orange!85!black,thick] (0,0) circle (6.5pt);
  \node[orange!85!black,fill=white,inner sep=1pt] at (0.15,0.42)
    {$\vec B_{\text{感}}$ 向外};
\end{tikzpicture}

\end{center}

\textbf{解题过程：}
\begin{enumerate}
  \item 原磁场垂直纸面向里，且磁感应强度增大，所以向里的磁通量增加。
  \item 感应磁场要阻碍这种增加，因此应垂直纸面向外。
  \item 右手拇指指向纸面外，四指弯曲方向为逆时针，所以从纸面正面观察，感应电流沿逆时针方向。
\end{enumerate}

\textbf{方法总结：} 不能从“磁场向里”直接猜电流方向。必须先判断磁通量是增加还是减少，再确定感应磁场方向。
\end{examplebox}

\begin{examplebox}{导体棒切割磁感线}
\textbf{题意分析：} 长 $L=0.40\ \mathrm{m}$ 的导体棒在 $B=0.50\ \mathrm{T}$ 的匀强磁场中以 $v=2.0\ \mathrm{m/s}$ 垂直切割磁感线。回路总电阻 $R=2.0\ \Omega$。求感应电动势和感应电流。

\textbf{物理模型：} 导体棒切割磁感线。先画出导轨、导体棒、磁场方向和速度方向，再判断能否直接用 $\mathcal E=BLv$。

\begin{center}
% TikZ 图：导体棒切割磁感线例题草图
% Day 23：电磁感应
\begin{tikzpicture}[scale=0.82,>=Latex,line cap=round]
  \path[use as bounding box] (-2.15,-1.35) rectangle (2.65,1.30);
  \draw[thick] (-1.85,0.85) -- (2.10,0.85);
  \draw[thick] (-1.85,-0.85) -- (2.10,-0.85);
  \draw[thick] (-1.85,0.85) -- (-1.85,-0.85);
  \draw[very thick,red] (-0.10,-0.85) -- (-0.10,0.85);
  \node[red] at (0.18,0.18) {$L$};
  \draw[->,red,very thick] (-0.10,0) -- (0.92,0) node[right] {$v$};
  \foreach \x in {-1.35,-0.75,-0.15,0.45,1.05,1.65} {
    \foreach \y in {-0.45,0.25} {
      \draw[blue,thick] (\x-0.055,\y-0.055) -- (\x+0.055,\y+0.055);
      \draw[blue,thick] (\x-0.055,\y+0.055) -- (\x+0.055,\y-0.055);
    }
  }
  \node[blue,font=\small] at (2.30,0.70) {$B$ 向里};
  \node[font=\small] at (0,-1.20)
    {先判定 $B,L,v$ 两两垂直，再用 $\mathcal E=BLv$};
\end{tikzpicture}

\end{center}

\textbf{解题过程：}
\[
\mathcal E=BLv=0.50\times0.40\times2.0=0.40\ \mathrm{V}
\]
\[
I=\frac{\mathcal E}{R}=\frac{0.40}{2.0}=0.20\ \mathrm{A}
\]

\textbf{方法总结：} 导体棒题先确认 $B$、$L$、$v$ 两两垂直，再用 $\mathcal E=BLv$；方向另用楞次定律或右手定则判断。
\end{examplebox}

\begin{examplebox}{线圈进出磁场的电流时段}
\textbf{题意分析：} 一个矩形线圈匀速进入一块匀强磁场区域，然后完全在磁场中运动，最后离开磁场。问哪些阶段有感应电流？

\textbf{物理模型：} 判断磁通量是否变化。不要只看“线圈在动”，要把进入、完全在内部、离开三个阶段分开画。

\begin{center}
% TikZ 图：线圈进出磁场例题草图
% Day 23：电磁感应
\begin{tikzpicture}[scale=0.78,>=Latex,line cap=round]
  \draw[thick,fill=blue!4] (-0.20,-0.85) rectangle (3.25,0.85);
  \foreach \x in {0.25,0.85,1.45,2.05,2.65} {
    \foreach \y in {-0.45,0.25} {
      \draw[blue,thick] (\x-0.05,\y-0.05) -- (\x+0.05,\y+0.05);
      \draw[blue,thick] (\x-0.05,\y+0.05) -- (\x+0.05,\y-0.05);
    }
  }

  \draw[thick,red] (-1.55,-0.42) rectangle (-0.70,0.42);
  \draw[->,red,thick] (-1.12,0) -- (-0.32,0);
  \node[font=\small,align=center] at (-1.12,-1.12) {进入\\$\Phi\uparrow$};

  \draw[thick,orange] (1.08,-0.42) rectangle (1.93,0.42);
  \draw[->,orange,thick] (1.35,0) -- (2.15,0);
  \node[font=\small,align=center] at (1.50,-1.12) {内部\\$\Phi$ 不变};

  \draw[thick,purple] (2.78,-0.42) rectangle (3.63,0.42);
  \draw[->,purple,thick] (3.05,0) -- (3.85,0);
  \node[font=\small,align=center] at (3.22,-1.12) {离开\\$\Phi\downarrow$};
\end{tikzpicture}

\end{center}

\textbf{解题过程：}
\begin{itemize}
  \item 进入磁场：穿过线圈的面积增大，磁通量增大，有感应电流。
  \item 完全在磁场中：面积不变、磁场不变，磁通量不变，无感应电流。
  \item 离开磁场：穿过线圈的面积减小，磁通量减小，有感应电流。
\end{itemize}

\textbf{方法总结：} 线圈题先看磁通量是否变化，不只看线圈是否运动；面积、磁场、夹角任一变化都可能产生感应电流。
\end{examplebox}

\begin{examplebox}{感应电荷量与变化快慢}
\textbf{题意分析：} 一个 $N=200$ 匝的闭合线圈总电阻为 $R=10\ \Omega$。每匝磁通量从 $5.0\times10^{-4}\ \mathrm{Wb}$ 单调减小到 $1.0\times10^{-4}\ \mathrm{Wb}$，忽略线圈自感。求这一过程中通过线圈某横截面的电荷量。

\textbf{物理模型：} 电阻不变、磁通量单调变化，感应电流方向不变，可以直接用累计电荷量公式。

\textbf{解题过程：}
\[
Q=\frac{N|\Delta\Phi|}{R}
=\frac{200\times(5.0-1.0)\times10^{-4}}{10}
=8.0\times10^{-3}\ \mathrm C.
\]

\textbf{方法总结：} 结果只由 $N$、$R$ 和总磁通量变化量决定，与完成变化所用时间无关；若磁通量来回变化导致电流反向，则应分段计算。
\end{examplebox}

\begin{examplebox}{发电机模型中的磁通量与感应电动势}
\textbf{题意分析：} $N=200$ 匝、面积 $S=2.0\times10^{-2}\ \mathrm{m^2}$ 的线圈在 $B=0.50\ \mathrm{T}$ 的匀强磁场中以角速度 $\omega=100\ \mathrm{rad/s}$ 匀速转动。取线圈面积法线与磁场方向重合的时刻为 $t=0$，转动方向使 $\theta=\omega t$ 增大，并按所选正法线用右手螺旋定则规定回路正方向。求感应电动势随时间的表达式、最大值及第一次达到正向最大值的时刻。

\textbf{物理模型：} 面积法线匀速转动，磁通量按余弦规律变化；感应电动势由磁通量的瞬时变化率决定。

\begin{center}
% TikZ 图：发电机磁通量与感应电动势的相位关系
% Day 23：电磁感应
\begin{tikzpicture}[scale=0.88,>=Latex,line cap=round,line join=round]
  \path[use as bounding box] (-0.55,-2.55) rectangle (6.65,2.35);

  % Magnetic flux: cosine.
  \draw[->,thick] (0,1.10) -- (6.15,1.10) node[right] {$t$};
  \draw[->,thick] (0,0.10) -- (0,2.18) node[above] {$\Phi$};
  \draw[blue,very thick,domain=0:360,samples=150]
    plot ({0.55+0.0145*\x},{1.10+0.82*cos(\x)});
  \node[blue,font=\small] at (4.15,2.05) {$\Phi=BS\cos\omega t$};

  % Induced emf: sine.
  \draw[->,thick] (0,-1.15) -- (6.15,-1.15) node[right] {$t$};
  \draw[->,thick] (0,-2.15) -- (0,-0.10) node[above] {$\mathcal E$};
  \draw[red,very thick,domain=0:360,samples=150]
    plot ({0.55+0.0145*\x},{-1.15+0.82*sin(\x)});
  \node[red,font=\small] at (4.98,-0.35)
    {$\mathcal E=\mathcal E_m\sin\omega t$};

  % Quarter-period guide.
  \draw[gray,dashed] (1.855,2.05) -- (1.855,-2.05);
  \fill[blue] (1.855,1.10) circle (1.8pt);
  \fill[red] (1.855,-0.33) circle (1.8pt);
  \node[font=\small,anchor=west] at (2.10,-2.30)
    {$t=T/4:\quad \Phi=0,\quad |\mathcal E|=\mathcal E_m$};
\end{tikzpicture}

\end{center}

\textbf{解题过程：}
\[
\Phi=BS\cos\omega t,
\qquad
\mathcal E=-N\frac{\mathrm d\Phi}{\mathrm dt}
=NBS\omega\sin\omega t.
\]
感应电动势最大值为
\[
\mathcal E_m=NBS\omega
=200\times0.50\times2.0\times10^{-2}\times100
=200\ \mathrm V.
\]
因此
\[
\boxed{\mathcal E=200\sin(100t)\ \mathrm V}.
\]
第一次达到正向最大值时 $\omega t=\pi/2$，所以
\[
\boxed{t=\frac{\pi}{2\omega}=\frac{\pi}{200}\ \mathrm s
\approx1.57\times10^{-2}\ \mathrm s}.
\]

\textbf{方法总结：} $t=0$ 时磁通量最大，但曲线斜率为零，所以感应电动势为零；经过四分之一周期，磁通量为零而变化率最大，感应电动势达到最大值。
\end{examplebox}

\section*{七、变式训练}

\begin{enumerate}
  \item 匀强磁场 $B=0.20\ \mathrm{T}$ 垂直穿过面积 $0.50\ \mathrm{m^2}$ 的线圈，求磁通量。
  \item 100 匝线圈中，每一匝的磁通量在 $0.10\ \mathrm{s}$ 内从 $0.020\ \mathrm{Wb}$ 变为 $0.005\ \mathrm{Wb}$，求平均感应电动势大小。
  \item 长 $0.30\ \mathrm{m}$ 的导体棒以 $4.0\ \mathrm{m/s}$ 垂直切割 $0.25\ \mathrm{T}$ 的磁场，求感应电动势。
  \item 一个线圈完全处在匀强磁场中匀速平移，是否一定有感应电流？说明理由。
  \item 一个总电阻为 $5.0\ \Omega$ 的单匝闭合线圈中，磁通量从 $0.060\ \mathrm{Wb}$ 单调减小到 $0.010\ \mathrm{Wb}$。若忽略自感且回路电阻保持不变，求整个过程中通过线圈某截面的电荷量。
  \item 闭合线圈中的磁场垂直纸面向外且正在减弱。从纸面正面观察，判断感应电流方向，并写出判断过程。
  \begin{center}
  \begin{tikzpicture}[scale=0.86,>=Stealth,line cap=round,line join=round]
  \path[use as bounding box] (-0.25,-0.92) rectangle (7.65,2.30);
  \begin{scope}
    \fill[blue!5] (-0.05,0.15) rectangle (3.0,2.1);
    \foreach \x in {0.35,0.85,1.35,1.85,2.35,2.85}
      \foreach \y in {0.35,0.85,1.35,1.85}
        \node[blue!75] at (\x,\y) {$\times$};
    \draw[thick] (0.20,0.35) -- (2.85,0.35);
    \draw[thick] (0.20,1.80) -- (2.85,1.80);
    \draw[very thick,red] (1.15,0.35) -- (1.15,1.8);
    \draw[->,red,thick] (1.15,1.05) -- (2.0,1.05) node[midway,above,font=\small] {$v$};
    \draw[<->,black] (1.02,0.35) -- (1.02,1.8) node[midway,left,font=\small] {$L$};
    \node[font=\small,align=center] at (1.48,-0.55)
      {导体棒切割磁感线\\$\mathcal E=BLv$};
  \end{scope}
  \begin{scope}[xshift=4.55cm]
    \fill[blue!5] (-0.15,0.15) rectangle (2.65,2.1);
    \foreach \x in {0.2,0.7,1.2,1.7,2.2}
      \foreach \y in {0.35,0.85,1.35,1.85}
        \node[blue!75] at (\x,\y) {$\times$};
    \draw[thick] (0.55,0.55) rectangle (1.85,1.65);
    \draw[->,red,thick] (1.2,1.1) -- (2.18,1.1) node[midway,above,font=\small] {$v$};
    \node[font=\small,align=center] at (1.20,-0.55)
      {线圈整体平移\\磁通量不变};
  \end{scope}
\end{tikzpicture}

  \end{center}
  \item 一台理想交流发电机中，线圈从磁通量最大的位置开始匀速转动。经过四分之一周期时，磁通量和感应电动势分别处于什么状态？
\end{enumerate}

\textbf{画图提示：} 第 3、4 题先画“磁场区域、运动方向、有效面积”。第 3 题看导体棒是否切割磁感线；第 4 题不要只看线圈在动，要看穿过线圈的磁通量是否变化。第 5 题先确认磁通量单调变化；第 6 题按楞次定律四步判断，不能直接猜电流方向。

\begin{learningbox}{变式训练参考结果}
\begin{enumerate}[itemsep=0.45em,topsep=0.3em,parsep=0pt]
  \item $\Phi=BS=0.10\ \mathrm{Wb}$。
  \item $\bar{\mathcal E}=N|\Delta\Phi|/\Delta t=15\ \mathrm V$。
  \item $\mathcal E=BLv=0.30\ \mathrm V$。
  \item 不一定。若线圈始终完全处在匀强磁场中，且面积和取向不变，则磁通量不变，没有感应电动势，也没有感应电流。
  \item 由 $Q=\int |I|\,\mathrm dt=|\Delta\Phi|/R$，得
  $Q=(0.060-0.010)/5.0=1.0\times10^{-2}\ \mathrm C$。
  \item 原磁场向外且磁通量减小，感应磁场应向外以阻碍减小；由右手螺旋定则，从纸面正面观察，感应电流沿逆时针方向。
  \item 经过四分之一周期时，面积法线与磁场垂直，磁通量为零；此时磁通量变化率的绝对值最大，感应电动势大小达到最大值。
\end{enumerate}
\end{learningbox}

\section*{八、动力学与能量视角}

\begin{competitionbox}
\textbf{电磁感应中的能量守恒}

在光滑水平导轨上，导体棒在外力作用下匀速运动。若导轨和导体棒的电阻均忽略、回路中的定值电阻为 $R$，并忽略回路自感，则外力不断做功，电路中产生焦耳热。电磁感应不是“凭空产生电能”，而是把机械能转化为电能，再转化为内能。

\begin{center}
% TikZ 图：电磁感应能量转化
% Day 23：电磁感应
\begin{tikzpicture}[scale=0.9,>=Latex]
  \draw[thick] (-2,1) -- (2.3,1);
  \draw[thick] (-2,-1) -- (2.3,-1);
  \draw[very thick,red] (0,-1) -- (0,1);
  \draw[->,red,thick] (0,0.25) -- (1.05,0.25) node[right] {$v$};
  \draw[->,purple,thick] (0,-0.25) -- (1.05,-0.25) node[right] {$F_{\text{外}}$};
  \draw[->,orange,thick] (0,0) -- (-1.05,0) node[left] {$F_A$};
  \foreach \x in {-1.5,-0.8,-0.1,0.6,1.3,2.0} {
    \foreach \y in {-0.55,0.25} {
      \draw[blue,thick] (\x-0.06,\y-0.06) -- (\x+0.06,\y+0.06);
      \draw[blue,thick] (\x-0.06,\y+0.06) -- (\x+0.06,\y-0.06);
    }
  }
  \node[blue] at (2.75,0.8) {$B$ 向里};
  \node at (0,-1.55) {$P_{\text{外}}=F_{\text{外}}v=I^2R$};
\end{tikzpicture}


\figcaption{导体棒匀速切割磁感线：外力功率等于电路热功率}\label{fig:induction-energy-flow}
\end{center}

若导体棒长 $L$，速度 $v$，回路总电阻 $R$，则
\[
\mathcal E=BLv,\qquad I=\frac{BLv}{R}
\]
安培力大小
\[
F_A=BIL=\frac{B^2L^2v}{R}
\]
若导体棒匀速运动，外力 $F_{\text{外}}=F_A$，外力功率
\[
P_{\text{外}}=F_{\text{外}}v=\frac{B^2L^2v^2}{R}
\]
电路热功率
\[
P_R=I^2R=\left(\frac{BLv}{R}\right)^2R=\frac{B^2L^2v^2}{R}
\]
两者相等，体现能量守恒。
\end{competitionbox}

\subsection*{1.\quad 楞次定律的能量解释}

楞次定律说“阻碍磁通量变化”，不是说永远阻碍运动。若磁通量的变化由相对运动引起，感应电流产生的安培力常表现为阻碍相对运动；若磁通量的变化由磁场强弱改变引起，阻碍的就是磁通量增大或减小本身。这个“阻碍”保证了能量不会凭空产生。

\subsection*{2.\quad 导体棒的终端速度}

若导体棒在光滑水平导轨上受恒定外力 $F$ 拉动，磁场、有效长度和回路总电阻保持不变，并忽略自感，则速度越大，感应电动势越大，电流越大，安培力也越大：
\[
v\uparrow \Rightarrow \mathcal E=BLv\uparrow \Rightarrow I=\frac{\mathcal E}{R}\uparrow
\Rightarrow F_A=BIL\uparrow.
\]
随着 $F_A$ 增大，导体棒的加速度逐渐减小；其速度最终趋近某一稳定值。该稳定状态满足 $F_A=F$，相应速度称为终端速度。

\begin{center}
% TikZ 图：导体棒终端速度
% Day 23：电磁感应
\begin{tikzpicture}[scale=0.9,>=Latex]
  \draw[thick] (-2.2,1.0) -- (2.6,1.0);
  \draw[thick] (-2.2,-1.0) -- (2.6,-1.0);
  \draw[very thick,red] (0,-1.0) -- (0,1.0);
  \foreach \x in {-1.7,-0.9,-0.1,0.7,1.5,2.3} {
    \foreach \y in {-0.55,0.25} {
      \draw[blue] (\x-0.05,\y-0.05) -- (\x+0.05,\y+0.05);
      \draw[blue] (\x-0.05,\y+0.05) -- (\x+0.05,\y-0.05);
    }
  }
  \draw[->,red,thick] (0,0.35) -- (1.1,0.35) node[right] {$v$};
  \draw[->,purple,thick] (0,-0.15) -- (1.1,-0.15) node[right] {$F$};
  \draw[->,orange,thick] (0,-0.55) -- (-1.0,-0.55) node[left] {$F_A$};
  \node[blue] at (2.95,0.85) {$B$ 向里};
  \node[font=\small,align=center] at (0,-1.65)
    {$v$ 增大 $\Rightarrow E$ 增大 $\Rightarrow I$ 增大 $\Rightarrow F_A$ 增大};
  \node[font=\small,align=center] at (0,-2.25)
    {最终匀速：$F=F_A$};
\end{tikzpicture}


\figcaption{导体棒终端速度：速度增大引起反向安培力增大}\label{fig:motional-terminal-speed}
\end{center}

\begin{examplebox}{导体棒的终端速度}
\textbf{题意分析：} 长为 $L$ 的导体棒在磁感应强度为 $B$ 的匀强磁场中沿光滑水平导轨运动，回路总电阻为 $R$。若水平外力恒为 $F$，忽略自感，求最终匀速速度。

\textbf{物理模型：} 速度趋近终端值时，合力趋近于零。

\textbf{解题过程：}
\[
\mathcal E=BLv,\qquad I=\frac{BLv}{R}.
\]
安培力
\[
F_A=BIL=\frac{B^2L^2v}{R}.
\]
达到稳定状态时
\[
F=F_A=\frac{B^2L^2v}{R},
\]
所以
\[
v=\frac{FR}{B^2L^2}.
\]

\textbf{方法总结：} 电磁感应综合题常有反馈链：速度决定电流，电流决定安培力，安培力又决定速度变化。
\end{examplebox}

\subsection*{3.\quad 自感、互感与涡流：变化磁通量的不同表现}

电磁感应不只发生在线圈进出磁场或导体棒切割磁感线时。若线圈自身电流发生变化，线圈中的磁通量也随之变化，于是会在线圈自身产生感应电动势，这叫\textbf{自感}。自感电动势仍遵循楞次定律，其作用是阻碍线圈中电流的变化，而不是阻止电流本身。

若一个线圈中的电流变化，使附近另一个线圈中的磁通量发生变化，从而在另一个线圈中产生感应电动势，这叫\textbf{互感}。变压器和无线充电都利用了互感：一个回路的电流变化，通过变化磁场影响另一个回路。

\begin{center}
% TikZ 图：互感
% Day 23：电磁感应
\begin{tikzpicture}[>=Latex,scale=0.95,every node/.style={font=\scriptsize}]
  \draw[rounded corners=3pt,fill=gray!12,draw=gray!60,thick] (0.7,0.45) rectangle (6.8,2.05);
  \node[gray!70!black] at (3.75,2.28) {共同铁芯/共同磁通};

  \foreach \x in {1.25,1.55,1.85,2.15,2.45} {
    \draw[blue,thick] (\x,0.50) .. controls +(0,0.55) and +(0,-0.55) .. (\x+0.18,2.00);
  }
  \node[blue] at (1.95,0.10) {原线圈};

  \foreach \x in {4.95,5.25,5.55,5.85,6.15} {
    \draw[red,thick] (\x,0.50) .. controls +(0,0.55) and +(0,-0.55) .. (\x+0.18,2.00);
  }
  \node[red] at (5.65,0.10) {副线圈};

  \draw[->,blue,very thick] (0.35,1.25) -- (1.05,1.25) node[midway,above] {$I_1$ 变化};
  \draw[->,red,very thick] (6.85,1.25) -- (7.55,1.25) node[midway,above] {$E_2$};

  \draw[->,orange,thick] (2.8,1.25) -- (4.55,1.25);
  \node[orange,above] at (3.68,1.25) {$\Delta\Phi$};

  \node[align=center] at (3.75,-0.55) {原线圈电流变化 $\rightarrow$ 磁通量变化 $\rightarrow$ 副线圈产生感应电动势};
\end{tikzpicture}


\figcaption{互感：一个线圈中的电流变化，可在另一个线圈中产生感应电动势}\label{fig:mutual-induction-coils}
\end{center}

当金属块或金属板处在随时间变化的磁场中，或相对非匀强磁场运动而使其内部闭合路径的磁通量发生变化时，金属内部会出现环形感应电流，叫\textbf{涡流}。金属板在匀强磁场中整体匀速平动并不必然产生涡流。涡流可以产生发热、电磁阻尼或电磁驱动。

\begin{center}
\begin{tikzpicture}[scale=0.92,>=Stealth]
  \begin{scope}
    \node[font=\small] at (0,1.65) {自感};
    \draw[thick] (-2.3,0) -- (-1.25,0);
    \draw[thick] (-1.25,0) -- (-0.90,0.35);
    \draw[thick] (-0.65,0) -- (0.0,0);
    \draw[decorate,decoration={coil,aspect=0.45,segment length=4pt,amplitude=4pt},thick] (0,0) -- (1.2,0);
    \draw[thick] (1.2,0) -- (2.25,0);
    \draw[->,red,very thick] (-2.1,0.35) -- (-1.15,0.35);
    \node[font=\scriptsize,red] at (-1.55,0.67) {$I$ 增大};
    \draw[->,blue,very thick] (1.70,-0.38) -- (0.70,-0.38);
    \node[font=\scriptsize,blue] at (1.35,-0.72) {感应电动势阻碍变化};
  \end{scope}

  \begin{scope}[xshift=5.4cm]
    \node[font=\small] at (0,1.65) {涡流与电磁阻尼};
    \draw[thick,fill=gray!10] (-1.1,-1.0) rectangle (1.1,1.0);
    \foreach \x in {-0.75,0,0.75}{
      \foreach \y in {-0.55,0.25}{
        \node[font=\large] at (\x,\y) {$\times$};
      }
    }
    \draw[red,thick,->] (-0.45,0.15) arc (140:-180:0.42);
    \draw[red,thick,->] (0.55,-0.25) arc (-40:280:0.42);
    \draw[->,very thick,orange!85!black] (-1.65,0) -- (-1.12,0);
    \node[font=\scriptsize,orange!85!black] at (-1.60,0.34) {金属板运动};
    \draw[->,very thick,blue] (1.45,0) -- (0.95,0);
    \node[font=\scriptsize,blue,anchor=west] at (1.66,0.34) {阻尼力};
  \end{scope}
\end{tikzpicture}


\figcaption{自感与涡流：变化电流和变化磁通量都会激发阻碍变化的感应效应}\label{fig:self-induction-eddy-current}
\end{center}

\textbf{应用视角：} 自感是日光灯镇流器、开关电源和电感元件的基础；互感是变压器和无线充电的重要基础；涡流用于电磁炉加热，也会造成变压器铁芯发热，所以工程上常用硅钢片叠片来减小涡流损耗。

\textbf{学习边界：} 高中阶段先会用“磁通量变化”和楞次定律解释这些现象，不要求定量计算自感系数、互感系数或复杂涡流分布。遇到应用题时，先说清哪一个磁通量在变化、感应效应阻碍什么变化即可。

\section*{九、易错点}

\begin{mistakebox}
\begin{enumerate}
  \item \textbf{认为有磁场就一定有感应电流：} 先看穿过回路的磁通量是否变化；若要形成持续电流，还要检查回路是否闭合。
  \item \textbf{角度用错：} 写 $\Phi=BS\cos\theta$ 前先画平面法线，$\theta$ 是磁场与法线的夹角。
  \item \textbf{楞次定律说成“阻碍运动”：} 先说清原磁通量怎样变；感应电流阻碍的是磁通量变化，有时才表现为阻碍相对运动。
  \item \textbf{完全进入匀强磁场后仍认为有电流：} 分段看有效面积；完全进入且磁通量不变时，没有感应电动势。
  \item \textbf{混淆感应电动势和感应电流：} 先判断是否有电动势，再看回路是否闭合；不闭合时没有持续电流。
  \item \textbf{直接套用 $\mathcal E=BLv$：} 先确认有效长度、速度和磁场方向两两垂直；否则取有效分量或回到法拉第定律。
  \item \textbf{把 $N|\Delta\Phi|/R$ 当成任意过程的累计电荷量：} 先确认 $R$ 不变并忽略自感；若感应电流反向，还要分段区分累计电荷量与净电荷量。
\end{enumerate}
\end{mistakebox}

\section*{十、今日小结}

\begin{enumerate}
  \item 选定正法线后，磁通量 $\Phi=\vec B\cdot\vec S=BS\cos\theta$ 是有正负的标量。
  \item 磁通量变化会产生感应电动势；闭合导电回路中才能形成沿回路的感应电流。
  \item 楞次定律判断方向：感应电流产生的磁场阻碍引起它的磁通量变化。
  \item 平均感应电动势由 $\Delta\Phi/\Delta t$ 决定，瞬时感应电动势由 $\mathrm d\Phi/\mathrm dt$ 决定。
  \item $R$ 不变且忽略自感时，$|Q_{\text{净}}|=N|\Delta\Phi|/R$；电流反向时要分段处理。
  \item 导体棒垂直切割磁感线常用 $\mathcal E=BLv$；一般情形要先取有效分量或回到法拉第定律，能量来源是外力做功。
  \item 自感、互感和涡流都是电磁感应的具体表现，方向判断仍遵循楞次定律。
\end{enumerate}

\textbf{一分钟复盘：} 遇到电磁感应题，先在图上圈出回路并检查是否闭合，再判断磁通量是变大还是变小；用楞次定律写一句“感应磁场要阻碍什么变化”，最后再用法拉第定律算大小。方向、大小和闭合条件要分开检查。
